%%%%%% -*- Coding: utf-8-unix; Mode: latex

\documentclass[polish]{aghengthesis}
%\documentclass[english]{aghengthesis} %dla pracy w języku angielskim. Uwaga, w przypadku strony tytułowej zmiana
% języka dotyczy tylko kolejności wersji językowych tytułu pracy.

\usepackage{tgtermes}
\usepackage[T1]{fontenc}
\usepackage{polski}
\usepackage[utf8]{inputenc}
\usepackage{url}
\usepackage{subfigure}
\usepackage{tabularx}
\usepackage{ragged2e}
\usepackage{booktabs}
\usepackage{multirow}
\usepackage{grffile}
\usepackage{indentfirst}
\usepackage{caption}
\usepackage{listings}
\usepackage[ruled,linesnumbered,lined]{algorithm2e}
\usepackage[bookmarks=false]{hyperref}

\hypersetup{
    colorlinks=true,
    allcolors=black
}

\usepackage[svgnames]{xcolor}
\usepackage{inconsolata}

\usepackage{csquotes}
\DeclareQuoteStyle[quotes]{polish}
  {\quotedblbase}
  {\textquotedblright}
  [0.05em]
  {\quotesinglbase}
  {\fixligatures\textquoteright}
\DeclareQuoteAlias[quotes]{polish}{polish}

\usepackage[nottoc]{tocbibind}

\usepackage[
style=numeric,
sorting=nyt,
isbn=false,
doi=true,
url=true,
backref=false,
backrefstyle=none,
maxnames=10,
giveninits=true,
abbreviate=true,
defernumbers=false,
backend=biber]{biblatex}
\addbibresource{references.bib}

\lstset{
    %language=Python, %% PHP, C, Java, etc.
    basicstyle=\ttfamily\footnotesize,
    backgroundcolor=\color{gray!5},
    commentstyle=\it\color{Green},
    keywordstyle=\color{Red},
    stringstyle=\color{Blue},
    numberstyle=\tiny\color{Black},    
    % morekeywords={TestKeyword},
    % mathescape=true,
    escapeinside=`',
    frame=single, %shadowbox, 
    tabsize=2,
    rulecolor=\color{black!30},
    title=\lstname,
    breaklines=true,
    breakatwhitespace=true,
    framextopmargin=2pt,
    framexbottommargin=2pt,
    extendedchars=false,
    captionpos=b,
    abovecaptionskip=5pt,
    keepspaces=true,            
    numbers=left,                    
    numbersep=5pt,                  
    showspaces=false,                
    showstringspaces=false,
    showtabs=false,
    tabsize=2
  }

\SetAlgorithmName{\LangAlgorithm}{\LangAlgorithmRef}{\LangListOfAlgorithms}
\newcommand{\listofalgorithmes}{\tocfile{\listalgorithmcfname}{loa}}

\renewcommand{\lstlistingname}{\LangListing}
\renewcommand\lstlistlistingname{\LangListOfListings}

\renewcommand{\lstlistoflistings}{\begingroup
\tocfile{\lstlistlistingname}{lol}
\endgroup}

% Definicje nowych rodzajów kolumn w tabeli
\newcolumntype{Y}{>{\small\centering\arraybackslash}X}
%\newcolumntype{b}{>{\hsize=1.6\hsize}Y}
%\newcolumntype{m}{>{\hsize=.6\hsize}Y}
%\newcolumntype{s}{>{\hsize=.4\hsize}Y}

\captionsetup[figure]{skip=5pt,position=bottom}
\captionsetup[table]{skip=5pt,position=top}

%%%%%%%%%%%%%%%%%%%%%%%%%%%%%%%%%%%%%%%%%%%%%%%%%%%%%%%%%%%%%%%%%%%%%%%%%%%%%%%
\author{Jan Kowalski, Jan Malinowski, Wojciech Kowalski}

\titlePL{Społecznościowy system wspomagający zarządzanie odtwarzaniem muzyki w obiektach usługowych i użyteczności publicznej}
\titleEN{A community system supporting the management of music playback in service and public facilities}

\fieldofstudy{Informatyka}

%\typeofstudies{Stacjonarne}

\supervisor{dr hab.\ inż.\ Krzysztof Iksiński, prof.\ AGH}

\date{\the\year}

%%%%%%%%%%%%%%%%%%%%%%%%%%%%%%%%%%%%%%%%%%%%%%%%%%%%%%%%%%%%%%%%%%%%%%%%%%%%%%%
\begin{document}

\maketitle
\tableofcontents

%%%%%%%%%%%%%%%%%%%%%%%%%%%%%%%%%%%%%%%%%%%%%%%%%%%%%%%%%%%%%%%%%%%%%%%%%%%%%%%
\chapter{\ChapterTitleProjectVision}
\section{Charakterystyka problemu}

Podczas nauki języka obcego łatwo natrafić na liczne pułapki lingwistyczne, do
szczególnie uciążliwych należy dobrze znane zjawisko jakim jest rozpoznawanie
fałszywych przyjaciół i interferencji językowych. Mózg człowieka podczas
przyswajania nowej wiedzy (w tym języka obcego) naturalnie dąży do optymalizacji
i szuka wzorców. Uczący się podświadomie opierają się na strukturach swojego
języka ojczystego, co ułatwia start, ale jednocześnie otwiera drogę do błędów.
To zjawisko nazywamy \textbf{transferem językowym} \cite{transfer_jezyowy}. Transfer
może być pozytywny, wtedy pomaga (gdy reguły w obu językach są zbieżne), ale problemem,
na którym skupia się aplikacja, jest \textbf{transfer negatywny}, prowadzący do
powielania błędnych schematów.

\section{Fałszywi przyjaciele i interferencje językowe}

Fałszywi przyjaciele tłumacza (w językoznawstwie nazywane często tautonimami) to
pary wyrazów w dwóch językach, które wykazują znaczne podobieństwo graficzne lub
fonetyczne, w rzeczywistości mają zupełnie inne znaczenia występujące w co
najmniej dwóch różnych językach.

Interferencje językowe to nie tylko słownictwo, ale też przenoszenie
struktur gramatycznych, składniowych czy nawet interpunkcyjnych. Przenikanie cech
jednego języka do drugiego, co zazwyczaj prowadzi do błędów w nauce języków obcych
lub w tłumaczeniach. Nazywa się je kalkami językowymi. Tłumacz próbuje "przetłumaczyć"
swoje myśli z języka ojczystego na język obcy słowo po słowie, co skutkuje zdaniami
niezrozumiałymi dla native speakera. Rozpoznawanie fałszywych przyjaciół jest w istocie
walką z jedną z najbardziej zdradliwych form interferencji.

Geneza zjawiska może pochodzić ze wspólnego pochodzenia etymologicznego. Wiele słów
w językach europejskich ma wspólne korzenie (np. łacińskie lub greckie), ale na
przestrzeni wieków ich znaczenie w poszczególnych krajach ewoluowało w zupełnie
innych kierunkach. W niektórych przypadkach może to być przypadkowa zbieżność brzmień.
Jest to rzadsza sytuacja, w której słowa brzmią identycznie, ale nie mają ze sobą nic
wspólnego historycznie.

Przykłady fałszywych przyjaciół dla pary języków polski-angielski, które pozwalają
lepiej zobrazować przedstawione zjawisko:
\begin{itemize}
    \item \textbf{lecture (ang.) / lektura (pol.)} - w języku angielskim oznacza
    wykład, na przykład na uczelni, a w języku polskim, podobne słowo, lektura
    oznacza książkę, którą studenci powinni przeczytać w ramach kursu lub sam proces
    czytania książki. Użycie tego słowa przez studenta szukającego "lektury"
    w zagranicznej bibliotece doprowadzi do nieporozumienia, gdyż zostanie on
    skierowany na wykład.
    \item \textbf{fatigue (ang.) / fatyga (pol.)} - choć oba słowa mają związek
    z wysiłkiem, ich zakres semantyczny jest inny. Fatigue to stan fizjologicznego
    wyczerpania, podczas gdy polska fatyga to raczej kłopot, trud podjęty dla kogoś.
    \item \textbf{actually (ang.) / aktualnie (pol.)} - bardzo częsty błąd, gdzie
    \textit{actually} oznacza "właściwie/w rzeczywistości", a polskie "aktualnie"
    to po angielsku \textit{currently}. Użycie tego słowa zmienia cały sens wypowiedzi
    o czasie.
\end{itemize}

\section{Motywacja projektu}

Podręczniki zazwyczaj wymieniają fałszywych przyjaciół w formie suchych tabel,
które są trudne do zapamiętania i nużące. Brakuje w nich interaktywności
i kontekstu. Aplikacja mobilna celuje w rozwiązanie tego problemu poprzez zmianę
paradygmatu nauki. Wykorzystanie koncepcji \textit{microlearningu} (nauki
w małych dawkach) oraz mechanizmów grywalizacji pozwala na budowanie skojarzeń
w sposób bezstresowy.

Zależy nam na wsparciu użytkowników w nauce języków obcych ze szczególnym
uwzględnieniem nauki fałszywych przyjaciół w szybki, zabawny i przyjemny sposób.
Chcielibyśmy zwrócić uwagę na istotność problemu fałszywych przyjaciół tłumacza,
które często występują między językami pokrewnymi.

\section{Istniejące rozwiązania}
Przed rozpoczęciem prac warto zrobić przegląd alternatywnych rozwiązań, już
dostępnych na rynku dla użytkownika.

TODO zdjęcia

\subsection{Duolingo}
Duolingo \cite{duolingo} to obecnie najpopularniejsza na świecie bezpłatna aplikacja i platforma
internetowa do nauki języków obcych, opierająca się na koncepcji mikronauki oraz
silnej grywalizacji (punkty doświadczenia, ligi, wirtualne waluty) w celu motywacji
i systematyczności użytkowników. Nauka przypomina zabawę i obejmuje krótkie,
interaktywne lekcje, ćwiczenia wymowy, czytania oraz tłumaczenia.

Mimo ogromnej popularności, Duolingo jest platformą ogólnego przeznaczenia.
Platforma nie posiada dedykowanego modułu skupiającego się na pułapkach
lingwistycznych, takich jak fałszywi przyjaciele czy interferencje językowe.
Użytkownik nie ma możliwości wyizolowania problematycznych par wyrazów (np.
lecture vs lektura) i przećwiczenia ich w skoncentrowany sposób. Ponadto aplikacja
rzadko tłumaczy kontekst błędu - po prostu wskazuje złą odpowiedź.

\subsection{Anki}
Anki \cite{anki} to darmowy program opierający się na cyfrowych fiszkach i systemie powtórek
w odstępach czasu (SRS). Algorytm optymalizuje czas nauki, wyświetlając słówka
tuż przed momentem, w którym mózg by je zapomniał. Umożliwia wykorzystywanie
gotowych talii, albo wgrywanie własnych fiszek.

Głównym problemem Anki jest surowość interfejsu oraz wymóg dużej ingerencji ze
strony użytkownika. Aby skutecznie uczyć się fałszywych przyjaciół, użytkownik
musi najpierw samodzielnie wyszukać słówka, stworzyć talię kart i zdefiniować
reguły. Ponadto w Anki brakuje elementów grywalizacji i różnorodności trybów nauki,
co może szybko znudzić użytkownika.

% \subsection{Ubiegłoroczna praca inżynierska}
% TODO

% \subsection{Podsumowanie}
% TODO

\section{Wizja produktu}
Zasadniczą wizją projektu jest stworzenie platformy edukacyjnej ukierunkowanej
na eliminację błędów wynikających z interferencji językowych oraz zjawiska
fałszywych przyjaciół tłumacza. Docelowym środowiskiem jest platforma mobilna,
co ma zapewnić użytkownikom możliwość nauki w dowolnym miejscu i czasie oraz
dostępność natychmiastowych powiadomień. Z uwagi na dostępność materiałów, system
w pierwszej fazie rozwoju zostanie zaimplementowany dla pary języków
polski-angielski z możliwością skalowalności.

\subsection{Aplikacja}
Aplikacja zaoferuje predefiniowane, wbudowane zestawy zadań podzielone na
kategorie tematyczne lub poziomy zaawansowania. Dodatkowo w przyszłości wdrożony
zostanie moduł importu, pozwalający użytkownikom na wgrywanie i tworzenie
własnych, spersonalizowanych list słówek. System umożliwi założenie konta, co
pozwoli na personalizację doświadczeń, gromadzenie statystyk oraz śledzenie
postępów w nauce.

\subsection{Spaced Repetition (Powtórki w interwałach)}
% TODO cytowania

Metoda \emph{Spaced Repetition} służy zapamiętywaniu informacji w wydajniejszy
sposób niż zwyczajne powtórki. Polega na odwlekaniu powtarzania dobrze wyuczonych
informacji na późniejszy termin i skupianiu się na świeżych informacjach.

Została po raz pierwszy opisana przez Arthura W. Meltona w jego pracy "The
Situation with Respect to the Spacing of Repetitions and Memory" z 1970
roku [cytowanie]. Tam odwoływał się do powszechnie znanego prawa
Jolt'a postulującego, że "jeżeli dwa ślady pamięci są odświeżone i przetestowane
w tej samej chwili, to w następującym okresie czasu nowszy ślad zaniknie prędzej
od starszego." [cytowanie] Melton wskazał, że tzw. "Total Time Law" (prawo
postulujące, że sumaryczny czas poświęcany na powtórki jest najważniejszy) się nie
sprawdza w praktyce i możliwe jest wydajniej zarządzać czasem nauki.

Istnieją różne algorytmy opierające się na zasadzie powtórek w interwałach. Jednym
z najpopularniejszych jest SM-2 stworzony na potrzeby oprogramowania SuperMemo
[cytowanie].

% TODO więcej o SM-2.

% TODO więcej o FSRS

\subsection{Mechanika ćwiczeń}
Główną częścią edukacyjną aplikacji będą interaktywne tryby ćwiczeń zaprojektowane
tak, aby zwiększać zainteresowanie użytkownika nauką. Wśród dostępnych rodzajów
ćwiczeń znajdą się:
\begin{itemize}
    \item \textbf{Testy jednokrotnego wyboru} - użytkownik będzie musiał wskazać
    poprawne znaczenie słowa spośród czterech dostępnych opcji. Kluczowym elementem
    tego zadania jest to, że jedną z błędnych odpowiedzi będzie zawsze znaczenie
    fałszywego przyjaciela.
    \item \textbf{Łączenie słów w pary} - zadanie polegające na połączeniu słowa
    w języku obcym z jego właściwym odpowiednikiem. Trudność będzie polegać na
    wymieszaniu par, w których fałszywi przyjaciele będą stanowić utrudnienie.
    \item \textbf{Analiza definicji} - bardziej zaawansowany tryb, w którym
    użytkownik dopasowuje słowo do jednej z czterech obcojęzycznych definicji.
    Podobnie jak w teście jednokrotnego wyboru, jedna z definicji będzie celowo
    opisywać fałszywego przyjaciela.
\end{itemize}

% \subsection{System grywalizacji}
% TODO

\section{Studium wykonalności}
\subsection{Architektura systemu}
System zostanie oparty na scentralizowanej architekturze klient-serwer. Wynika
to bezpośrednio z założeń wizji produktu: konieczności utrzymywania globalnych
tabel rankingowych, synchronizacji profili użytkowników oraz przechowywania
jednolitego zbioru danych z ćwiczeniami. Rozdzielenie warstwy klienckiej od logiki
biznesowej i składowania danych pozwoli w przyszłości na łatwe dodanie kolejnych
klientów, np. aplikacji webowej bez konieczności przepisywania całego systemu.

\subsection{Warstwa serwerowa}
Warstwa serwerowa zostanie zaimplementowana w języku Python. Język ten
charakteryzuje się wysoką czytelnością i zwięzłą składnią. Pozwoli to na
przejrzyste prezentowanie kluczowych algorytmów, np. logiki przydzielania punktów
czy wybierania zestawów zadań, a także łatwość implementacji i wyszukiwania
błędów. Warto wspomnieć, że Python posiada wiele użytecznych frameworków, np.
\texttt{FastAPI}, które pozwalają na wystawienie interfejsu API niezbędnego do
komunikacji z aplikacją mobilną.

\subsection{Warstwa kliencka}
Ze względu na powszechność użycia, obsługę natychmiastowych powiadomień, oraz budowę
nawyku codziennej nauki, jako docelową platformę kliencką wybrano system
\textbf{Android}. Aplikacja zostanie napisana w nowoczesnym, oficjalnie zalecanym
przez Google, języku \textbf{Kotlin}. Wykorzystanie nowoczesnych, deklaratywnych
narzędziami do budowy interfejsu takich jak \textbf{Jetpack Compose}, zagwarantuje
płynność działania i łatwość utrzymania kodu.

Ustalono minimalną wspieraną wersję na \textbf{Android 8.0} (API 26), co zapewnia
pokrycie większości urządzeń na rynku. Aplikacja będzie kompilowana pod docelową
wersję \textbf{Android 16}, co pozwoli na wykorzystanie najnowszych funkcjonalnosci.

% \subsection{Baza danych}
% TODO

\subsection{API}
Komunikacja pomiędzy warstwą kliencką a serwerową będzie realizowana za pomocą
bezstanowego interfejsu \textbf{REST API} z wykorzystaniem protokołu
\textbf{HTTP}, co zapewni łatwość w zarządzaniu, skalowalność i możliwość łatwego
rozwinięcia projektu o kolejne systemy. Wymiana danych będzie odbywać się przy
użyciu formatu \textbf{JSON}. Wybór struktury JSON jest uzasadniony jej lekkością,
wysoką czytelnością oraz niskim narzutem sieciowym.

\section{Analiza zagrożeń}
W trakcie realizacji projektu mogą wystąpić różnego rodzaju zagrożenia zarówno
związane z kwestiami technicznymi, jak i koncepcyjnymi. Poniżej przedstawiono
omówienie głównych czynników ryzyka.
\begin{enumerate}
    \item \textbf{Zbiór danych} - problemem może być pozyskanie odpowiedniej ilości
    danych pozwalających na utworzenie przykładów do nauki wybranego języka, które
    pozwolą na efektywną naukę.

    \textbf{Proponowane rozwiązanie:} \\
    Na początkowym etapie projektu uwaga zostanie skupiona na stworzeniu
    ograniczonego, ale jakościowego zbioru danych dla pary języków
    polski-angielski. W przyszłości baza może zostać rozbudowana poprzez
    integrację z zewnętrznymi API słownikowymi, wykorzystanie metod ekstrakcji
    danych z otwartych źródeł lub wdrożenie mechanizmu pozwalającego społeczności
    na dodawanie własnych zestawów.
    \item \textbf{Zainteresowanie użytkownika} - wystarczająca rożnorodność gier
    i poziom zainteresowania użytkownika, ponieważ bardzo łatwo jest, w tego typu
    aplikacjach, wywołać efekt znudzenia u użytkownika.
    
    \textbf{Proponowane rozwiązanie:} \\
    Istotne jest żeby utrzymać użytkownika w aplikacji przez odpowiednią długość
    czasu co pozwoli poprawić efekty nauki. Zastosowanie mechanizmów
    grywalizacji, takich jak ograniczenia czasowe, system punktacji oraz tabele
    rankingowe. Ponadto, architektura aplikacji zakłada modułowość, co pozwoli na
    zmiany trybów nauki w ramach jednej sesji, zapobiegając w ten sposób zjawisku
    monotonii.
    \item \textbf{Trudność zadań} - dobranie odpowiedniego poziomu trudności zadań
    językowych tak, żeby uzyskać wystarczający poziom nauki, ale nie zniechęcić za
    szybko użytkownika trudnością zadań.

    \textbf{Proponowane rozwiązanie:} \\
    Podział bazy słów na poziomy zaawansowania. Dodatkowo zaimplementowany
    zostanie algorytm adaptacyjny, który będzie analizował historię odpowiedzi
    użytkownika i na jej podstawie dostosowywał odpowiedni zbiór danych.
    \item \textbf{Przyjazny interfejs} - ryzyko stworzenia nieintuicyjnego lub
    mało estetycznego interfejsu graficznego, co negatywnie wpłynie na odbiór
    aplikacji.

    \textbf{Proponowane rozwiązanie:} \\
    Zastosowanie popularnych schematów projektowania przejrzystych interfejsów
    oraz wykorzystanie w projekcie deklaratywnej biblioteki interfejsu Jetpack
    Compose znacznie ułatwi zachowanie spójności wizualnej, płynnych animacji
    i responsywności aplikacji.
\end{enumerate}

% \section{Słownik pojęć}
% TODO

%%%%%%%%%%%%%%%%%%%%%%%%%%%%%%%%%%%%%%%%%%%%%%%%%%%%%%%%%%%%%%%%%%%%%%%%%%%%%%%
\chapter{\ChapterTitleScope}
\label{sec:zakres-funkcjonalnosci}

% poniższą zawartość rodziału należy usunąć z finalnej wersji pracy.
\emph{Kontekst użytkowania produktu (aktorzy, współpracujące systemy) oraz specyfikacja wymagań funkcjonalnych i niefunkcjonalnych.}

%%%%%%%%%%%%%%%%%%%%%%%%%%%%%%%%%%%%%%%%%%%%%%%%%%%%%%%%%%%%%%%%%%%%%%%%%%%%%%%
\chapter{\ChapterTitleRealizationAspects}
\label{sec:wybrane-aspekty-realizacji}

% poniższą zawartość rodziału należy usunąć z finalnej wersji pracy.
\emph{Przyjęte założenia, struktura i zasada działania systemu, wykorzystane rozwiązania technologiczne wraz z uzasadnieniem ich wyboru, istotne mechanizmy i zastosowane algorytmy.}

%%%%%%%%%%%%%%%%%%%%%%%%%%%%%%%%%%%%%%%%%%%%%%%%%%%%%%%%%%%%%%%%%%%%%%%%%%%%%%%
\section{Algorytmy}
\label{sec:algorytmy}

\begin{algorithm}[!htbp]
\SetKwData{Left}{left}\SetKwData{This}{this}\SetKwData{Up}{up}
\SetKwFunction{Union}{Union}\SetKwFunction{FindCompress}{FindCompress}
\SetKwInOut{Input}{input}\SetKwInOut{Output}{output}
\Input{A bitmap $Im$ of size $w\times l$}
\Output{A partition of the bitmap}
\BlankLine
\emph{special treatment of the first line}\;
\For{$i\leftarrow 2$ \KwTo $l$}{
\emph{special treatment of the first element of line $i$}\;
\For{$j\leftarrow 2$ \KwTo $w$}{\label{forins}
\Left$\leftarrow$ \FindCompress{$Im[i,j-1]$}\;
\Up$\leftarrow$ \FindCompress{$Im[i-1,]$}\;
\This$\leftarrow$ \FindCompress{$Im[i,j]$}\;
\If(\tcp*[h]{O(\Left,\This)==1}){\Left compatible with \This}{\label{lt}
\lIf{\Left $<$ \This}{\Union{\Left,\This}}
\lElse{\Union{\This,\Left}}
}
\If(\tcp*[f]{O(\Up,\This)==1}){\Up compatible with \This}{\label{ut}
\lIf{\Up $<$ \This}{\Union{\Up,\This}}
\tcp{\This is put under \Up to keep tree as flat as possible}\label{cmt}
\lElse{\Union{\This,\Up}}\tcp*[h]{\This linked to \Up}\label{lelse}
}
}
\lForEach{element $e$ of the line $i$}{\FindCompress{p}}
}
  \caption[Przykładowy algorytm]{Przykładowy algorytm.}
  \label{alg:pseudo-code}
\end{algorithm}

%%%%%%%%%%%%%%%%%%%%%%%%%%%%%%%%%%%%%%%%%%%%%%%%%%%%%%%%%%%%%%%%%%%%%%%%%%%%%%% 
\section{Fragmenty kodu źródłowego}
\label{sec:listingi}
Listing~\ref{lst:maximum} przedstawia przykładowy fragment kodu źródłowego.

\begin{lstlisting}[language=Python,float=!htbp,
    caption={[Przykładowy fragment kodu]Przykładowy fragment kodu },label=lst:maximum]
# The maximum of two numbers

def maximum(x, y):

    if x >= y:
        return x
    else:
        return y

x = 2
y = 6
print(maximum(x, y),"is the largest of the numbers ", x, " and ", y)

\end{lstlisting}

%%%%%%%%%%%%%%%%%%%%%%%%%%%%%%%%%%%%%%%%%%%%%%%%%%%%%%%%%%%%%%%%%%%%%%%%%%%%%%%
\chapter{\ChapterTitleWorkOrganization}
\label{sec:organizacja-pracy}

% poniższą zawartość rodziału należy usunąć z finalnej wersji pracy.
\emph{Struktura zespołu (role poszczególnych osób), krótki opis i uzasadnienie przyjętej metodyki i/lub kolejności prac, planowane i zrealizowane etapy prac ze wskazaniem udziału poszczególnych członków zespołu, wykorzystane praktyki i narzędzia w zarządzaniu projektem.}

%%%%%%%%%%%%%%%%%%%%%%%%%%%%%%%%%%%%%%%%%%%%%%%%%%%%%%%%%%%%%%%%%%%%%%%%%%%%%%%
\chapter{\ChapterTitleResults}
\label{sec:wyniki-projektu}

% poniższą zawartość rodziału należy usunąć z finalnej wersji pracy.
\emph{Wskazanie wyników projektu (co konkretnie udało się uzyskać: oprogramowanie, dokumentacja, raporty z testów/wdrożenia, itd.), prezentacja wyników i ocena ich użyteczności (jak zostało to zweryfikowane --- np.\ wnioski klienta/użytkownika, zrealizowane testy wydajnościowe, itd.), istniejące ograniczenia i propozycje dalszych prac.}

%%%%%%%%%%%%%%%%%%%%%%%%%%%%%%%%%%%%%%%%%%%%%%%%%%%%%%%%%%%%%%%%%%%%%%%%%%%%%%%
\printbibliography
\listoffigures
\listoftables
\listofalgorithmes
\lstlistoflistings

\end{document}
